\chapter{Kammerets Dagligdag}
Selvom \TKET er en fest- og foredragsforening, sker der mange andre ting omkring foreningen. Siden 1967 har foreningen haft et lokale på B-gangen\footnote{Ny Munkegade 118, Bygning 1535, lokale 128}, hvor mange hænger ud og holder sin dagligdag på studiet. Derfor kommer et kapitel om hvordan dagligdagen er på \TKET.

%lokaler
\section{\TKET}
Da Matematisk Institut flyttede ind i dens nuværende bygninger på Ny Munkegade, var \TKET selvfølgelig blevet skrevet ind i byggeplanen. Derfor flyttede foreningen ind i det nu berømte lokale, der i dagligtale bliver kaldt: \TKET eller bare Kammeret. Selvom lokalet ikke er meget større end 4,3 $\times$ 4,7 meter, så er Kammeret fyldt med ting og sager. 

Langs den ene side stilles alle kasserne med drikkevare, som kommes i køleskabet, der står i det ene hjørne. Modsat der er tavlen, hvorpå \textit{Omgangsskyldnere}-feltet og $\Delta I\Sigma KOK\Upsilon\Gamma\Lambda EN$
er på. Lige ved siden af tavlen hænger klokken, som når man ringer med, kommer en omgang (én kasse) på bordet i midten af lokalet. 

På væggene rundt i hele lokalet, hænger mange forskellige skilte, billeder, og andet. Ved døren hænger \textit{Papmache-koen}, der er en udstoppet ged. Ved vinduet, hænger en diskokugle (ikke til forveksling med $\Delta I\Sigma KOK\Upsilon\Gamma\Lambda EN$) og True- og False skiltene. En af de mere kendte skilte er selvfølgelig \textit{Ingen-Sigma}-skiltet. Og da reglen er: ``Skilte på \TKET er regler", må man ikke skrive de store summeringstegn inde i lokalet.

En stor del af dagligdagen på Kammeret er Bestyrelsens \textit{Kontortid}. Mellem 11.15 og 12.00 i hverdagen, dukker et bestyrelsesmedlem op, rydder op, og laver én hel kande ``kontortids"-kaffe. Dette tidspunkt er også en tid, som mange benytter sig af at spise sin frokost i lokalet. På samme måde bruger mange også \TKET til at studere. Selvom man nemt kan blive distraheret af spændende samtaler, så er det også et sted, hvor man kan få hjælp til opgaver og TØ. Man vil jo hellere lave en andens opgave, end ens egen.

Selvom folk fordyber sig i deres studie på Kammeret, så er alting ikke lige så seriøst. Især efter studerende er færdig med de sidste forelæsninger kl 16, er en fyraftensøl ofte velkommen. Og mange dage kan en enkel fyraftensøl ofte smage af én til. Så bliver den tidlige aften et tidspunkt med hygge, spil og de bedste samtaler om alt mellem himmel og flaske-bunden.

\section{Hængerne}
Som du nok læser i skriftet, bliver ordet \textit{hængere} ofte brugt, her vil vi forklare hvad det egentlig er. Definitionen for hængere har ændret sig med tiden, men de fleste er enige om at det er en person der kommer semi ofte på \TKET{}, måske de hjælper bestyrelsen til deres kontortid, eller andre daglige ting, måske de hjælper med at stå i bar til en \TK{}-fest. Altså humlen er at de kommer på \TKET og hjælper gerne foreningen. 
%	TK
\section{Te-Køkkenet}
Lige ved siden af \TKET på B-gangen, ligger \textit{The-Køkkenet}. Her er udover en vask og service, er der også en mikrobølgeovn, normal ovn, og 4 kogeplader. Her kan man varme sin frokost og fylde sin drikkedunk. Men ofte vil køkken-delen blive brugt af bestyrelsen og andre til at lave aftensmad.

\section{Rådet}
%	Rådet
På den anden side af gangen, ligger Rådet. Et studie- og mødelokale varetaget af Fagrådet Sigma. Dog bruger \TKET og Mat/Fys-tutorforening også lokalet til at have møder. Ellers bruges det til at side og studere. På en torsdag aften efter et torsdagsmøde, sætter folk sig også ofte derind for at spille spil, og snakke.

\section{Facebook grupper}
I takt med at verden er blevet mere digitaliseret, er der også kommet Facebook grupper til at dele forskellige informationer/vittigheder. Den største og "vigtigeste" er Kammerets Daglidag. Med godt og vel 700 medlemmer, er det en nem måde at række ud til folk på. Her laves opslag om alt mulig information, om det er happenings til en forelæsning, nogen der laver aftensmad en dag i Te-køkkenet, eller om DR der vil dele en dokumentar de har lavet, hvor \TKET er med.

\subsection{TK-Memes}
I en læseferie i 2019, lavede en hænger nogle sjove memes og delte med sine venner, dette var optakten til den næststørste Facebook gruppe for folk på \TKET{}. Her deles som navnet nok siger, memes om alt fra en sjov ting der skete til et torsdagsmøde, eller noget sjovt der skete til et foredrag.

%Kontortid og dagligdag
%Aften
